% Background and Motivation: concepts, limitations, evidence, and challenges.
\section{Background and Motivation}
\label{sec:background-motivation}

\subsection{Background}
\label{sec:motivation-background}

\paragraph{Multi-agent execution.}
A multi-agent application divides a user task among agents that exchange
intermediate results. Each agent maintains a conversation and advances
through LLM requests and tool calls: the model generates a response or
a tool invocation, and the returned result becomes input to a subsequent
request~\cite{wu2023autogen}. A \emph{request} is one model invocation;
an \emph{agent} can issue multiple requests to complete its task; a
\emph{workflow} connects these tasks through sequential dependencies,
parallel branches and joins. The application orchestrates this workflow,
while the serving engine executes its model requests.

Figure~\ref{fig:agent-workflow-examples} illustrates this structure.
For code generation, a main agent delegates a CSV-export feature to API,
UI, test-authoring and documentation agents. Starting from a shared
specification, each branch alternates model requests with file reads,
edits or checks. The main agent then integrates the changes and runs
tests. For deep research, a lead agent delegates technical, cost, safety
and supply investigations, then synthesizes their evidence into a
report~\cite{anthropic2025multiagent}. Branches can require different
numbers of requests and tools; there is no common round barrier.
A completed model request can therefore leave both its agent and the
workflow unfinished.

\paragraph{Workflow completion time.}
Let $w$ identify a workflow, $a_w$ denote its arrival time at the system,
and $c_w$ denote the time when all its required outputs are ready.
We define \emph{workflow completion time} (WCT) as the elapsed time from
arrival to completion, $\mathrm{WCT}_w=c_w-a_w$.
It includes model execution, tool execution,
queueing and dependency waits. In the coding example, integration waits
for all required changes. A delay on the last unfinished branch postpones
integration, while an earlier-finishing branch may absorb the same delay
before the join. Request latency affects WCT through these dependencies.

\begin{figure}[t]
  \centering
  \includegraphics[width=\columnwidth]{figures/agent-workflow-examples.pdf}
  \caption{Illustrative multi-agent workflows.
    Each row is one agent: $R_i$ denotes its $i$th LLM request, with
    intervening tool results carried into the next request.
    Branches have different request counts and no shared round barriers.
    The coordinator resumes after their results join.
    Positions and lengths do not encode execution time.}
  \label{fig:agent-workflow-examples}
  \Description{Two stacked left-to-right workflows. A main agent assigns
    CSV-export API, UI, test-authoring and documentation tasks in parallel,
    then integrates their outputs and runs tests. Their request counts
    are four, three, three and two, respectively. A lead researcher assigns
    technical, cost, safety and supply investigations, with four, two,
    three and two requests, before writing a report. Each branch explicitly
    unfolds its requests and intervening tools. Starts and ends are
    staggered; all edges move forward, without loops or round barriers.}
\end{figure}

\paragraph{KV reuse in multi-agent workflows.}
During \emph{prefill}, an LLM processes input tokens and computes
per-layer keys and values (KV); \emph{decoding} extends this state as
output tokens are generated. With compatible model and position semantics,
requests with an identical token prefix can reuse its KV and compute
only the remaining suffix~\cite{zheng2024sglang}.

Multi-agent workflows create reuse both within and across agents.
Within an agent, successive requests often append model outputs and
tool results to the existing conversation, allowing earlier turns'
KV to be reused~\cite{gao2024cachedattention}. Across agents, branches
can inherit a common prompt before adding their own tasks. For example,
the API, UI and test agents in Figure~\ref{fig:agent-workflow-examples}
can share KV for common project instructions and repository context
when these form the same initial token sequence. Reuse extends only
through the matching prefix; different role prompts or context edits
can shorten it~\cite{zheng2024sglang}.

Retained KV can therefore serve an agent's later turns or another
agent's requests, including after a tool call or dependency wait.
KV consists of model-generated tensors rather than the token IDs stored
by the application. In conventional per-token layouts, its size grows
with context length, layer count and KV-head dimensions. Retaining
long histories can thus consume substantially more memory than saving
their text or token IDs.

\paragraph{Failures and KV reconstruction.}
Failures are a recurring concern in large-scale LLM serving deployments:
hardware faults and software crashes can interrupt workers and make
their GPU-resident KV unavailable~\cite{cao2026lumen,jayakody2026ghostserve}.
Redirecting an interrupted request to a worker with the model loaded
does not restore this state. Even when the application retains the
conversation's tokens, the model must reconstruct the missing KV before
it can resume generation. Without a usable KV copy, recovery requires
processing the retained history again~\cite{strati2024dejavu}.

This reconstruction is particularly costly for agents that accumulate
long contexts over many requests~\cite{jayakody2026ghostserve}.
Coding agents retain source files, edits and test outputs; research
agents collect retrieved documents, intermediate findings and tool
results. A failure late in execution can therefore require rebuilding
KV for the accumulated history, not just the latest user input.
In a multi-agent workflow, one failed worker may interrupt several
co-located agents, leaving the surviving workers to reconstruct multiple
histories while also handling normal traffic~\cite{cao2026lumen}.
Reconstruction delays the affected branches and can postpone downstream
integration or synthesis, increasing workflow completion time.

\paragraph{KV retention across storage tiers.}
A \emph{KV checkpoint} saves an already computed prefix. If that copy
survives a failure, recovery can retrieve it and recompute only the
remaining tail~\cite{strati2024dejavu,jayakody2026ghostserve}.
Hierarchical caches extend GPU HBM with host DRAM and local or shared
storage to retain more histories~\cite{gao2024cachedattention,qin2025mooncake}.
Shared storage, such as a shared filesystem, makes retained KV accessible
to multiple serving workers. A compatible worker can retrieve this state
after reassignment or failover, provided the data and storage service
survive the failure. Sharing describes the access scope, not a specific
medium such as SSD.

Location determines the access path: KV already in the executor's HBM
is directly available to kernels; local DRAM requires host-to-device
transfer, and remote DRAM also requires network transfer. With host
staging, shared storage adds storage and network reads before the
host-to-device transfer. Computation and prefetching can overlap some
of this work. Retention nevertheless occupies
capacity, and copying or moving KV consumes bandwidth that may also serve
foreground requests. Discarding the last reusable KV copy releases space
but requires reconstruction if that state is needed again.

\paragraph{Existing mechanisms.}
Hierarchical caching reduces repeated prefix work across
turns~\cite{gao2024cachedattention,qin2025mooncake}; KV streaming and
checkpointing reduce reconstruction after
failures~\cite{strati2024dejavu,jayakody2026ghostserve}.
Parrot exposes application dependencies for serving optimization, while
KVFlow uses upcoming agent steps to guide eviction and
prefetching~\cite{lin2024parrot,pan2025kvflow}. LUMEN coordinates
checkpoint placement and recovery routing with worker
load~\cite{cao2026lumen}. These capabilities provide a basis for
workflow-aware recovery. Allocating finite protection resources also
requires addressing two limitations of the decision criteria.

\paragraph{Limitation 1: Request-level priorities incompletely express
workflow benefit.}
Recency, proximity to the next use and reconstruction cost describe when
state may be needed and how expensive it is to recover. Taken alone,
they do not determine how restoring a request sooner advances the
workflow's required outputs. An expensive history can belong to a branch
with time to spare, while a shorter recovery can delay the branch that
gates completion. Protection priorities must account for dependencies
and the branches' remaining progress.

\paragraph{Limitation 2: Recovery speed alone omits retention value and
resource costs.}
A branch that can tolerate recovery delay may still reuse its history
on the next normal request. Removing its last KV copy forfeits that reuse;
moving the copy to a slower tier preserves it but adds retrieval work.
Conversely, keeping every recoverable history in fast memory consumes
space needed by other histories and foreground execution. Minimizing
recovery latency in isolation therefore does not resolve which state
to retain, where to retain it, or how much protection the serving system
can afford.

\begin{figure*}[t]
  \centering
  \captionsetup[subfigure]{font=normalsize,labelfont=normalfont,
    textfont=normalfont,justification=centering,singlelinecheck=false}
  \begin{subfigure}[t]{2.48in}
    \vspace{0pt}
    \begin{minipage}[c][1.65in][c]{\linewidth}
      \centering
      \includegraphics[width=\linewidth]{figures/workflow.pdf}
    \end{minipage}
    \caption{Four-branch workflow}
    \label{fig:motivation-workflow}
  \end{subfigure}\hfill
  \begin{subfigure}[t]{1.86in}
    \vspace{0pt}
    \begin{minipage}[c][1.65in][c]{\linewidth}
      \centering
      \includegraphics[width=\linewidth]{figures/workflow-gaps.pdf}
    \end{minipage}
    \caption{Workflow-maximum gaps}
    \label{fig:motivation-branch-gap}
  \end{subfigure}\hfill
  \begin{subfigure}[t]{2.52in}
    \vspace{0pt}
    \begin{minipage}[c][1.65in][c]{\linewidth}
      \centering
      \includegraphics[width=\linewidth]{figures/recovery-cost.pdf}
    \end{minipage}
    \caption{Context and Prefill cost}
    \label{fig:motivation-recompute}
  \end{subfigure}
  \caption{Workflow structure and the cost of reusable KV state.
    (a) Four trace-derived branches with schematic Main stages A0/A5.
    Bars preserve timing; dashed tails assume an all-branch join.
    (b) ECDF of maximum branch-duration gaps per workflow: 121 workflows
    with temporally inferred fan-in; these are not failure-time slack.
    (c) Measured Qwen2.5-7B-Instruct-1M Prefill on two A100 PCIe 40\,GB GPUs
    (TP2, BF16, batch one), with sparse attention above 32K.
    Reuse starts from $(C-1)$ cached tokens and computes one;
    DRAM/shared-storage I/O is serialized with computation.
    HBM/DRAM: median and min--max of three timings;
    shared storage/full prefill: one timing per length.
    Protocol: Appendix~\ref{sec:motivation-protocol}.
    K~=~1,024; M~=~1,000,000.}

  \label{fig:motivation}
  \Description{Left: a four-branch workflow with A2 completing last at
    151.6 seconds. Middle: an ECDF of workflow-maximum branch-duration
    differences with a median of 138 seconds and a 90th percentile of
    1173 seconds. Right: four measured Prefill curves for Qwen2.5-7B-Instruct-1M
    from 8K to 1,010,000 tokens: blue circles for HBM reuse, orange squares
    for serialized DRAM reuse, purple diamonds for serialized shared-storage
    reuse, and red triangles for full prefill. Both axes are logarithmic.}
\end{figure*}

\subsection{Motivation and Challenges}

We examine these limitations using agent execution traces and measured
KV readiness costs. The traces characterize branch timing; the model
measurements quantify reconstruction and reuse under a common endpoint.
Together they motivate two requirements for workflow-aware protection.

\paragraph{Unequal progress across parallel branches.}
\label{sec:motivation-workflow}
Figure~\ref{fig:motivation}(\subref{fig:motivation-workflow}) shows four
trace-derived branches. A2 finishes last, 151.6\,s after the fork;
A1 finishes at 60.4\,s. If the following Main stage needs all four outputs,
A1 can absorb 91.2\,s of additional delay before postponing the join,
whereas any delay on A2 postpones it. We call the additional delay that
a branch can absorb without postponing the required output its
\emph{slack}.

The public SemiAnalysis CC/Weka trace contains 393 filtered Claude Code
sessions~\cite{semianalysis2026cctraces}. We retain overlapping subagent
groups whose branches complete and are followed by a Main request after
the last branch finishes; Main activity during the branches is allowed.
This temporal fan-in proxy retains 286 groups in 121 workflows.
For each group, we compute the longest minus the shortest branch duration,
then take each workflow's maximum gap.
Figure~\ref{fig:motivation}(\subref{fig:motivation-branch-gap}) shows a
median of 138\,s and a 90th percentile of 1,173\,s. These statistics show
substantial duration differences within the retained workflows.
The trace does not record explicit result-consumption edges, so they
neither verify all-branch joins nor measure slack at a fault. The
four-branch example is one retained group, not a representative
workflow-maximum gap.

\paragraph{From reconstruction time to workflow delay.}
To quantify the consequence of a join, we measure full prefill at the
four active requests' input-size proxies. Both this experiment and the
tier comparison in Figure~\ref{fig:motivation}(\subref{fig:motivation-recompute})
use Qwen2.5-7B-Instruct-1M~\cite{qwen2025qwen25oneM} on two A100 PCIe
40\,GB GPUs, with TP2, BF16 and batch one. Our instrumented SGLang
0.5.10rc0 execution uses MInference vertical/slash
attention~\cite{jiang2024minference} and dual chunk attention (DCA)
position handling, with sparse attention above 32K.
The four history-length proxies come from trace hash-block counts times
64 and range from 100,864 to 113,664 tokens. Each measured input is a
fixed synthetic history plus one token. After one complete warmup per
length, we collect three timings and report their median and min--max.
Appendix~\ref{sec:motivation-protocol} gives the complete protocol.

\begin{figure}[t]
  \centering
  \includegraphics[width=\linewidth]{figures/recovery-exposure.pdf}
  \caption{Measured Prefill costs for
    Figure~\ref{fig:motivation}(\subref{fig:motivation-workflow}), using
    Figure~\ref{fig:motivation}(\subref{fig:motivation-recompute})'s configuration
    at trace input-size proxies plus one token.
    Bars/whiskers show median/min--max over three timings after warmup.
    Hatching denotes absorbed delay; solid fill is the calculated WCT
    increase under an all-branch join with other execution fixed.
    Each case is a single-branch counterfactual, not a fault-injected
    WCT measurement.}
  \label{fig:motivation-recovery-exposure}
  \Description{A1 through A4 have measured median Prefill costs of
    33.2, 32.6, 35.4 and 39.9 seconds.
    Applying each cost separately as an added branch pause gives WCT
    increases of 0.0, 32.6, 16.3 and 18.0 seconds.
    Hatching identifies the absorbed portion; whiskers show the minimum
    and maximum of three timings, not confidence intervals.}
\end{figure}

Figure~\ref{fig:motivation-recovery-exposure} maps these costs to the
assumed all-branch join. At 29.5\,s after the fork, each case separately
adds a measured reconstruction pause $R_i$ to branch $i$, holding all
other execution and queues fixed. With offline slack $\sigma_i$ from
the trace completion times, the WCT increment is
\[
  \Delta\mathrm{WCT}_i=\max(0,R_i-\sigma_i).
\]
A1's 33.2\,s reconstruction fits within its 91.2\,s slack and leaves WCT
unchanged. A2's slightly smaller cost, 32.6\,s, is fully exposed because
its slack is zero. A3 and A4 expose 16.3 and 18.0\,s, respectively.
Thus similar reconstruction costs can have different workflow effects.
The costs are physical measurements; the WCT increments are conditional
calculations using synthetic inputs and assumed dependencies, not
fault-injected measurements of the original workflows.

\paragraph{Challenge 1: Estimating the workflow impact of recovery.}
A protection decision must connect the delay it avoids to the required
workflow outputs. The example uses completed traces to obtain slack,
but an online system sees only execution so far: remaining model and
tool durations are uncertain, and later dependencies may not yet be
exposed. Priorities must change as branches progress. A worker failure
can also interrupt several branches whose recoveries share resources,
so their isolated WCT increments cannot simply be added. The challenge
is to derive useful recovery priorities from partial workflow progress
while accounting for these uncertainties and interactions.

\paragraph{The cost of retaining and accessing long histories.}
\label{sec:motivation-long-context}
Figure~\ref{fig:motivation}(\subref{fig:motivation-recompute}) compares
reconstruction with reuse from HBM, DRAM and shared storage under the
same execution configuration. For $C=M+1$ tokens, full prefill starts
from empty KV; reuse starts from the same $M$-token history and computes
one new token. All timers end when the $C$ tokens' KV is ready, before
vocabulary projection. HBM is a resident-state reference; DRAM and shared
storage serialize restoration with computation. These measurements
exclude checkpoint creation, failure detection and recovery routing.

At 1,010,000 tokens, full prefill takes 571.21\,s despite sparse attention.
HBM, DRAM and shared-storage reuse take 1.305, 2.663 and 10.156\,s,
respectively. Retaining the history avoids most reconstruction work;
its location determines the additional retrieval delay. Shared storage
adds 7.49\,s over DRAM in this run, while the history occupies 57.92\,GB
across the two ranks. This separates the value of retaining reusable
state from the additional value of making it quickly accessible.
HBM/DRAM points report three timings, whereas full-prefill and
shared-storage points each report one. They characterize this fixed
input and serialized implementation; storage caches and I/O overlap can
change the differences. They do not measure the net serving benefit
of a protection policy.

\paragraph{Challenge 2: Balancing fast recovery, reuse and resource use.}
Finite capacity prevents every long history from remaining near
execution. Protection must decide both which histories retain reusable
state and which state merits faster access. These decisions
need to combine the workflow impact of recovery with future normal reuse,
capacity occupancy and data movement. In particular, low recovery
urgency does not imply low reuse value, and a faster restore path is
useful only to the extent that it advances required work. The challenge
is to allocate protection resources as these demands change without
allowing protection itself to undermine foreground service.

\label{sec:motivation-goals}
AegisKV targets these two challenges: using workflow progress to guide
recovery priorities, and managing limited KV capacity across storage
tiers while preserving normal reuse. The next section develops the
observations and mechanisms that support these decisions.
